\subsection{Proof of the implication ``MAT-free \texorpdfstring{$\Rightarrow$}{implies} strongly chordal''}
\label{subsec:MAT-free implies strongly chordal}
 
In this subsection we prove the ``only if'' part of our main Theorem \ref{thm:MAT-free-strong-chordal} (MAT-freeness implies strong chordality).

First we need a few preliminary results.
Let $ G = (V_{G}, E_{G}) $ be a simple graph. 
Let~$ \chi_G(t)$ be the \textbf{chromatic polynomial} of $G$ (the polynomial that counts the number of proper vertex colorings of $G$). 
It is known that $ \chi_G(t) = \chi_{ \mathcal{A}_{G} }(t) $ (\eg  \cite[Theorem 2.88]{OT92}).
Let $ \omega(G) $ denote the \textbf{clique number}, the number of vertices in a largest clique of $ G $. 
Recall that $\mcK(G)$ denotes the set of all maximal cliques of $G$.

\begin{proposition}\label{the number of maximal exponents = the number of largest cliques}
Let $ G $ be a chordal graph. Then the following statements hold. 
\begin{enumerate}[(1)]
\item The maximal exponents of $ \mathcal{A}_{G} $ are equal to $ \omega(G)-1 $. 
In addition, the number of maximal exponents of $ \mathcal{A}_{G} $ equals the number of largest cliques of $ G $.  
\item  If $ \lambda $ is an MAT-labeling of  $ G $, then the endvertices of each $ e \in  \pi_{n}$ where $n = \omega(G)-1$ are contained in a unique maximal clique of $G$. 
Furthermore, the map $ \phi : \pi_{n} \to \mcK(G)$ defined by $\phi(e)=$ the maximal clique containing the endvertices of $ e$ induces a bijection $ \pi_{n} \simeq \phi(\pi_{n})$ and $\phi(\pi_{n})=\{\mbox{all largest cliques of $ G $}\}$. 
\end{enumerate}
 
\end{proposition}
\begin{proof}
First we prove (1).
When $ G = K_{\ell} $, $ \omega(G) = \ell $ and $ \chi_{G}(t) = t(t-1) \cdots (t-\ell+1) $. 
The assertions clearly hold. 

We may assume that $ \ell \geq 3 $ and $ G $ is not complete. 
We proceed by induction on~$ \ell = |V_{G}| $. 
Since $ G $ is not complete, there exist two nonadjacent vertices $ a,b \in V_{G} $. 
Let $ S \subseteq V_{G} $ be a minimal $ (a,b) $-separator. 
Then $ V_{G} $ is decomposed as $ V_{G} = A \sqcup S \sqcup B $, where $ a \in A $ and $ b \in B $. 
Let $ G_{1} \coloneqq G[A\sqcup S] $ and $ G_{2} \coloneqq G[S \sqcup B] $. 
Note that these~$ G_{1} $ and $ G_{2} $ are chordal. 
Moreover, $ G[S] $ is complete (Theorem \ref{Dirac minimal vertex separator}) and $ \chi_{G}(t) = \chi_{G_{1}}(t)\chi_{G_{2}}(t)/\chi_{G[S]}(t) $ (\eg  \cite[Theorem 3]{Read68}). 

Let $ m_{1} $ and $ m_{2} $ denote the numbers of cliques consisting of $ \omega(G) $ many vertices in $ G_{1} $ and $ G_{2} $, respectively. 
Since there is no clique of $ G $ containing both vertices in~$ A $ and in $ B $, the number of largest cliques of $ G $ equals $ m_{1} $ + $ m_{2} $. 
By the induction hypothesis, the chromatic polynomials of $ G_{1} $ and $ G_{2} $ can be expressed as $ \chi_{G_{1}}(t)  = (t-\omega(G)+1)^{m_{1}} f(t) $ and $ \chi_{G_{2}}(t) = (t-\omega(G)+1)^{m_{2}}g(t) $, where $ f(t), g(t) \in \mathbb{Z}[t] $ are the products of some linear factors with roots strictly smaller than $ \omega(G) -1$. 
Therefore $ \chi_{G}(t) = (t-\omega(G)+1)^{m_{1}+m_{2}}f(t)g(t)/\chi_{G[S]}(t) $. 
Since $ \chi_{G[S]}(t) = t(t-1) \cdots (t-|S|+1) $ and $ |S| < \omega(G) $ (Remark \ref{Ho and Lee}), the maximal exponents of $ G $ are equal to $ \omega(G)-1 $  and the number of maximal exponents of $ G $ is $ m_{1} + m_{2} $. 

Now we prove part (2). If $ n = 0 $, the assertions hold trivially. 
If $ n = 1 $, then $ G $ is a forest and $ \lambda $ is a constant labeling whose value is $ 1 $. 
Hence the assertions also hold. 

Now suppose $ n \geq 2 $. 
Let $ e \in \pi_{n} $. 
Clearly, there exists $ C\in \mcK(G)$ such that $e \in G[C]$.
Suppose that there exist two distinct $ C_{1} ,C_{2} \in \mcK(G)$ such that $ e\in G[C_{1} \cap C_{2}] $. 
Since~$ \lambda_{E_{G \setminus e}} $ is an MAT-labeling, $ \mathcal{A}_{G\setminus e} $ is free and hence $ G\setminus e $ is chordal. 
By the maximality of $ C_{1} ,C_{2}$, there exist $u\in C_{1}\setminus C_{2} $ and $v\in C_{2}\setminus C_{1} $ such that $\{u,v\} \notin E_G$.
Then $u,v$ and the endvertices of $ e$ form a $4$-cycle  which is chordless in $ G\setminus e $, a contradiction. 
Thus the endvertices of each edge in $ \pi_{n} $ are contained in exactly one maximal clique of $ G $.
Therefore the map $\phi$ is well-defined. Moreover, $ |\phi(\pi_{n})| \leq  |\pi_{n}|  $. 

Let $ G^{\prime} \coloneqq G\setminus \pi_{n} $. 
The restriction $ \lambda|_{E_{G^{\prime}}} $ is an MAT-labeling of $G'$ hence $ \mathcal{A}_{G^{\prime}} $ is MAT-free. 
Therefore the maximal exponents of $ \mathcal{A}_{G^{\prime}} $ are equal to $ n-1 $.
By part (1), $ \omega(G^{\prime}) = \omega(G)-1 $. 
Thus every largest clique of $ G $ contains the endvertices of at least one edge in $ \pi_{n} $. 
Hence every largest clique of $ G $ belongs to $ \phi(\pi_{n}) $ and $ |\phi(\pi_{n})| \geq  |\pi_{n}|  $. 
Thus $ \phi(\pi_{n}) $ is precisely the set consisting of the largest cliques. 
This completes the proof. 
\end{proof}

In general, a restriction of an MAT-labeling is not an MAT-labeling (\eg  when we restrict to any edge in $\pi_k$ with $k>1$). 
We show below that it is the case for restriction to certain subset. 
Recall that  $ \mathcal{P}_{G} $ denotes the clique intersection poset of a chordal graph $ G $ (\S \ref{sec:more}). 
\begin{lemma}\label{restriction to intersection of maximal cliques}
If $ \lambda $ is an MAT-labeling of a chordal graph $ G $, then the restriction $ \lambda|_{E_{G[X]}} $ is an MAT-labeling of the subgraph $ (V_{G}, E_{G[X]}) $ for any node $ X \in \mathcal{P}_{G} $. 
\end{lemma}
\begin{proof}
We proceed by induction on $ n = \omega(G)-1 $. 
Again it is easily seen that the assertion holds true for  $ n = 0,1 $.
Now suppose $ n \geq 2 $. 
First we consider the case $X=C$ where $C$ is a largest clique of $ G $. 
Note that by Proposition \ref{restriction to subset}, it suffices to prove that $ \lambda|_{E_{G[C]}} $ satisfies \ref{definition MAT-labeling triangle}. 

Let $e=\{u,v\} \in\pi_n$ be the unique edge in $\pi_n$ whose endvertices are contained in~$C$ (Proposition \ref{the number of maximal exponents = the number of largest cliques}).
The clique $ C $ is the union of two largest cliques $ C^{\prime} = C \setminus  \{v\}$ and $ C^{\prime\prime} = C \setminus  \{u\} $ of $ G^{\prime} := G\setminus\pi_{n} $. 
For $ W \in \{C,C^{\prime}, C^{\prime\prime}\} $, define $ \pi_{k}^{W} \coloneqq \lambda|_{E_{G[W]}}^{-1}(k) = \pi_{k} \cap E_{G[W]} $ for $ k \in [n]$. 
Note that $ \pi_{n}^{C} = \{e\} $ and $ \pi_{k}^{C} = \pi_{k}^{C^{\prime}} \cup \pi_{k}^{C^{\prime\prime}} $ for $ k \in [n-1] $ since the restrictions $ \lambda|_{E_{G[C^{\prime}]}} $ and $ \lambda|_{E_{G[C^{\prime\prime}]}} $ are MAT-labelings by the induction hypothesis. 

To show \ref{definition MAT-labeling triangle} of $ \lambda|_{E_{G[C]}} $, it suffices to prove that every edge in $ \pi_{k}^{C} $ forms at least~$ k-1 $ triangles with edges in $ E_{k-1} \cap E_{G[C]} $ because every edge in $ \pi_{k}^{C} $ forms at most $ k-1 $ triangles with edges in $ E_{k-1} \cap E_{G[C]} $ by \ref{definition MAT-labeling triangle} of $ \lambda $. 
When $ k = n $, the edge~$ e \in \pi_{n}^{C} $ forms $ n-1 $ triangles with edges in $E_{G[C]}\setminus\{e\} \subseteq E_{n-1} \cap E_{G[C]} $. 
We are left with~$ k \in [n-1] $. 
Let~$ f \in \pi_{k}^{C} $, then $ f \in \pi_{k}^{C^{\prime}} \cup \pi_{k}^{C^{\prime\prime}} $. 
Without loss of generality, we may assume~$ f \in \pi_{k}^{C^{\prime}} $. 
Then $f$ forms $ k-1 $ triangles with edges in~$ E_{k-1} \cap E_{G[C^{\prime}]} \subseteq E_{k-1} \cap E_{G[C]} $. 
Therefore $ \lambda|_{E_{G[C]}} $ satisfies \ref{definition MAT-labeling triangle} and hence $ \lambda|_{E_{G[C]}} $ is an MAT-labeling. 

Now we treat the case $ X \in \mathcal{P}_{G} $ is not a largest clique of $ G $. 
First we claim that~$ E_{G[X]} \cap \pi_{n} = \varnothing $. 
If not, we can find an edge $ e \in E_{G[X]} \cap \pi_{n} $. 
Let $ C $ denote the largest clique such that $ e \in E_{G[C]} $. 
By the definition of $X$, there exists a maximal clique $ D $ of $G$ such that $X\subseteq  D \neq C$. 
Thus $ e \in E_{G[C\cap D]} $. 
Take a vertex $ c \in C\setminus D $. 
By the maximality of $ D $, there exists $ d \in D $ such that $\{c,d\} \notin E_G$. 
Then we obtain a chordless $4$-cycle  in $ G^{\prime} = G\setminus \pi_{n} $ formed by $ c,d $ and the endvertices of $ e $, which contradicts to the chordality of $ G^{\prime} $. 
Thus $ E_{G[X]} \cap \pi_{n} = \varnothing $. 

 


Observe that any maximal but not a largest clique of $ G $ is a maximal clique of $ G^{\prime} $. 
If $X$ is the intersection of non-largest maximal cliques in $G$, then $ X \in \mathcal{P}_{G^{\prime}} $.
By the induction hypothesis, $ \lambda|_{E_{G[X]}} = \left(\lambda|_{E_{G^{\prime}}}\right)|_{E_{G^{\prime}[X]}} $ is an MAT-labeling.
Otherwise, $X$ is contained in a largest clique of $G$, say $C$.
Let $e=\{u,v\} \in\pi_n$ be the unique edge in~$\pi_n$ whose endvertices are contained in $C$. 
 Then $ C = C^{\prime} \cup C^{\prime\prime}  $ where $ C^{\prime} = C \setminus  \{v\}$ and~$ C^{\prime\prime} = C \setminus  \{u\} $ are largest cliques of $ G^{\prime} $. 
 By the preceding paragraph, $e \notin E_{G[X]}$ hence either $u \notin X$ or $v \notin X$. 
 Thus either $X \subseteq C''$ or $X \subseteq C'$.
Hence we may replace~$C$ by $C'$ or $C''$ in the definition of $X$.
Therefore the node $ X $ is the intersection of some maximal cliques of $ G^{\prime} $. 
Again the  induction hypothesis applies. 
\end{proof}



We are ready to prove the main result of this subsection.
\begin{theorem}[MAT-freeness implies strong chordality]
\label{thm:main-1} 
If a simple graph $G$ admits an MAT-labeling, then $G$ is strongly chordal.  
